\documentclass[11pt,a4paper]{article}
\usepackage[UTF8]{ctex}
\usepackage{amsmath}
\usepackage{booktabs}
\usepackage{geometry}
\usepackage{xcolor}
\usepackage{hyperref}

\geometry{left=2.5cm,right=2.5cm,top=2.5cm,bottom=2.5cm}

\title{\textbf{Task 2.1: 构建词典与语料库技术报告}}
\date{2026年2月3日}

\begin{document}

\maketitle

\section{任务概述}

Task 2.1的目标是为每个时间窗口构建高质量的Gensim词典和BOW语料库，通过频率过滤去除噪声词和停用词，确保LDA模型训练在高质量词汇基础上进行。本任务采用三重过滤策略：最小文档频率、最大文档频率比例和词汇量上限。

\section{词汇过滤策略}

\subsection{过滤参数}

\begin{table}[h]
\centering
\begin{tabular}{ll}
\toprule
\textbf{参数} & \textbf{值} \\
\midrule
最小文档频率 (min\_df) & 100 \\
最大文档频率比例 (max\_df) & 0.4 \\
最大保留词汇量 (keep\_n) & 30,000 \\
\bottomrule
\end{tabular}
\caption{词汇过滤参数配置}
\end{table}

\subsection{过滤原理}

\begin{itemize}
    \item \textbf{min\_df=100}：删除在少于100个文档中出现的罕见词
    \item \textbf{max\_df=0.4}：删除在超过40\%文档中出现的通用词（停用词）
    \item \textbf{keep\_n=30,000}：保留最频繁的30,000个词
\end{itemize}

\textcolor{blue}{\textbf{设计理念}}：平衡词汇质量与计算效率，避免噪声词影响主题学习，同时控制模型复杂度。

\section{处理结果汇总}

\subsection{各窗口统计}

\begin{table}[h]
\centering
\begin{tabular}{lrrrrrr}
\toprule
\textbf{时间窗口} & \textbf{文档数} & \textbf{原始词汇} & \textbf{过滤后词汇} & \textbf{保留率} & \textbf{平均词数/文档} \\
\midrule
2016-2017 & 1,635,960 & 197,800 & 15,411 & 7.8\% & 17.1 \\
2018-2019 & 496,041 & 112,525 & 7,907 & 7.0\% & 16.5 \\
2020-2021 & 1,448,911 & 234,258 & 22,021 & 9.4\% & 27.3 \\
2022-2023 & 4,188,301 & 383,201 & 30,000 & 7.8\% & 26.5 \\
2024-2025 & 1,390,688 & 218,383 & 19,796 & 9.1\% & 28.3 \\
\midrule
\textbf{合计} & 9,160,136 & - & 95,135 & \textbf{8.3\%} & 23.8 \\
\bottomrule
\end{tabular}
\caption{词典构建结果汇总}
\end{table}

\subsection{质量指标分析}

\begin{itemize}
    \item \textbf{平均保留率}：8.3\%（标准差0.8\%）
    \item \textbf{词汇规模}：7,907 - 30,000词（取决于窗口大小）
    \item \textbf{文档覆盖}：99.99\%文档非空（过滤后）
    \item \textbf{平均文档长度}：16.5 - 28.3词/文档
\end{itemize}

\section{词汇分布分析}

\subsection{高频词谱系（Top 30）}

保留的高频词主要反映职业岗位的核心要求：

\begin{table}[h]
\centering
\begin{tabular}{llr}
\toprule
\textbf{排名} & \textbf{词汇} & \textbf{文档频率} \\
\midrule
1 & 管理 & 993,898 \\
2 & 大专 & 941,186 \\
3 & 熟练 & 762,498 \\
4 & 责任心 & 714,843 \\
5 & 本科 & 672,130 \\
6 & 协调 & 663,987 \\
7 & 行业 & 649,171 \\
8 & 学习 & 543,669 \\
9 & 销售 & 516,060 \\
10 & 项目 & 445,019 \\
\midrule
21-30 & 客户, 抗压, 善于, 培训, 技术, 开发, 团队合作\_精神, 业务 & 263K-310K \\
\bottomrule
\end{tabular}
\caption{保留高频词分布（前10名+21-30名）}
\end{table}

\textcolor{blue}{\textbf{洞察}}：高频词谱系反映了岗位招聘的核心维度：
\begin{itemize}
    \item 教育背景：大专、本科
    \item 技能要求：熟练、学习、技术、开发
    \item 个人素质：责任心、协调、抗压、善于
    \item 业务领域：管理、销售、项目、客户、业务
\end{itemize}

\subsection{中频词示例}

中频词代表专业领域术语：

\begin{verbatim}
服务性 (566), 商务部 (566), 清查 (566), 副校长 (565),
搅拌 (565), ELISA (565), 元至 (565), 定岗_定编 (565),
台面 (565), 一份_稳定_收入 (565), 离岗 (565), 电子仪器 (565),
优秀学生_干部 (565), 可能_出现 (565), Android_SDK (565)
\end{verbatim}

\subsection{低频词示例}

低频词反映特定行业或新兴领域：

\begin{verbatim}
过道 (132), CODE (132), 营建 (131), 有亲 (131),
底新 (131), 房等 (131), 工效 (131), 投诉率 (131),
二外 (131), 小吃街 (131), 广告网 (131), 榜单 (131),
加工资 (131), 直供 (131), 假牙 (131), 废旧物资 (131),
人好 (131), 单员 (131), life (131), 压裂 (131)
\end{verbatim}

\section{过滤效果评估}

\subsection{噪声去除分析}

\begin{table}[h]
\centering
\begin{tabular}{lrrr}
\toprule
\textbf{过滤类型} & \textbf{2016-2017} & \textbf{2022-2023} & \textbf{效果} \\
\midrule
原始词汇 & 197,800 & 383,201 & - \\
低频词移除 & 182,389 & 368,800 & 7.6\%-7.8\% \\
高频词移除 & 15,411 & 30,000 & 7.8\%-7.8\% \\
\midrule
最终保留 & 15,411 & 30,000 & 7.8\% \\
\bottomrule
\end{tabular}
\caption{词汇过滤效果对比}
\end{table}

\subsection{质量提升指标}

\begin{itemize}
    \item \textbf{词汇纯度}：移除通用停用词，提高主题区分度
    \item \textbf{计算效率}：控制词汇量在30,000以内，降低LDA复杂度
    \item \textbf{语义密度}：平均文档长度23.8词，信息密度适中
    \item \textbf{覆盖完整性}：保留率8.3\%，涵盖主要职业术语
\end{itemize}

\section{语料库构建}

\subsection{BOW格式转换}

\begin{itemize}
    \item \textbf{输入}：分词后的文档列表（每文档平均23.8词）
    \item \textbf{输出}：Matrix Market格式(.mm) + 索引文件(.mm.index)
    \item \textbf{存储效率}：稀疏矩阵格式，节省磁盘空间
    \item \textbf{加载速度}：Gensim原生支持，快速加载到内存
\end{itemize}

\subsection{语料统计}

\begin{table}[h]
\centering
\begin{tabular}{lrrr}
\toprule
\textbf{时间窗口} & \textbf{总词数} & \textbf{非空文档} & \textbf{平均长度} \\
\midrule
2016-2017 & 28,037,426 & 1,635,910 & 17.1 \\
2018-2019 & 8,205,478 & 495,998 & 16.5 \\
2020-2021 & 39,543,838 & 1,448,895 & 27.3 \\
2022-2023 & 111,000,468 & 4,188,267 & 26.5 \\
2024-2025 & 39,366,060 & 1,390,576 & 28.3 \\
\midrule
\textbf{合计} & 226,153,270 & 9,159,646 & \textbf{23.8} \\
\bottomrule
\end{tabular}
\caption{语料库统计信息}
\end{table}

\section{技术性能分析}

\subsection{处理效率}

\begin{itemize}
    \item \textbf{构建时间}：< 5分钟/窗口（取决于文档数量）
    \item \textbf{内存占用}：约2-8GB（取决于词汇量）
    \item \textbf{磁盘空间}：词典文件~400KB-780KB，语料库~10-50MB
\end{itemize}

\subsection{数据质量保证}

\begin{itemize}
    \item \textbf{完整性检查}：99.99\%文档成功转换
    \item \textbf{一致性验证}：词典与语料库词汇表匹配
    \item \textbf{可重现性}：相同参数保证结果一致
\end{itemize}

\section{输出文件}

\subsection{文件结构}

\begin{itemize}
    \item \texttt{window\_*\_dict}：Gensim词典文件（词汇表+频率统计）
    \item \texttt{window\_*\_corpus.mm}：BOW语料库（Matrix Market格式）
    \item \texttt{window\_*\_corpus.mm.index}：语料库索引文件
    \item \texttt{dictionary\_stats.pkl}：构建统计数据
\end{itemize}

\subsection{文件大小}

\begin{table}[h]
\centering
\begin{tabular}{lrrr}
\toprule
\textbf{窗口} & \textbf{词典大小} & \textbf{语料库大小} & \textbf{索引大小} \\
\midrule
2016-2017 & 403KB & 22MB & 3.2MB \\
2018-2019 & 207KB & 6.3MB & 0.8MB \\
2020-2021 & 581KB & 31MB & 4.6MB \\
2022-2023 & 780KB & 84MB & 12.6MB \\
2024-2025 & 516KB & 28MB & 4.2MB \\
\bottomrule
\end{tabular}
\caption{输出文件大小统计}
\end{table}

\section{结论}

Task 2.1成功构建了高质量的词典和语料库，关键成果包括：

\begin{itemize}
    \item ✅ \textbf{有效过滤}：平均保留8.3\%高质量词汇，去除噪声和停用词
    \item ✅ \textbf{语义丰富}：保留的核心词汇涵盖教育、技能、素质、业务四大维度
    \item ✅ \textbf{计算优化}：控制词汇量在30,000以内，保证LDA训练效率
    \item ✅ \textbf{数据完整}：99.99\%文档成功处理，总计2.26亿词元
    \item ✅ \textbf{格式标准}：输出Gensim兼容格式，便于后续LDA训练
\end{itemize}

该词典构建为LDA模型提供了优质的词汇基础，确保主题学习的准确性和可解释性，为后续岗位综合化分析奠定了坚实的数据基础。

\end{document}